\section{Motorcycles, mopeds and microcars}

\subsection{Motorcycles and mopeds until 2006}

Until 2006, motorcycles up to 50~cm³ and mopeds had a plate unlike a car's, issued
by the municipal council. Plates had the sequence \plate{0-XXX / 00-00}, where
\plate{XXX} is the three-letter code of the municipality where the vehicle was
registered (see appendix~\ref{sec:municipios}). This rule was set by
Regulatory Decree 13/98.

Moped plates had a yellow background and plates for motorcycles under 50~cm³ a
white one. Until the 1980s, yellow plates were also used on bicycles and
animal-drawn vehicles.

\platepair{moped-municipal-oer}{Moped from Oeiras: \plate{2-OER 20-75}.}%
          {moto-municipal-str}{Motorcycle up to 50~cm³ from Santarém: \plate{1-STR 00-30}.}

Decree-Law 106/2006 required these plates to be replaced by 2008, although some
stayed in use until 2010.

\subsection{Motorcycles and mopeds after 2006}

Motorcycles over 50~cm³ have plates with the same sequence as cars, on smaller
plates with no euroband and no yellow band showing the date of first
registration.

Motorcycles under 50~cm³ and mopeds have the same plate format as motorcycles
over 50~cm³, but with a yellow background.

\platepair{moto-6060uj}{Motorcycle: \plate{60-60-UJ}.}%
          {moped-08ns44}{Moped: \plate{08-NS-44}.}

\platefig[0.3\linewidth]{moto-49bx10}{Motorcycle: \plate{49-BX-10}.}

\subsection{Microcars (the “papa-reformas”)}

They have the usual plate layout, but with a yellow background instead of white.
They are four-wheeled mopeds (quadricycles): those with room for it may carry a
car-format plate with a yellow background instead of the motorcycle plate.

\platefig{minicar-77mq39}{Microcar: \plate{77-MQ-39}.}
